% Part: first-order-logic
% Chapter: introduction
% Section: models-theories

\documentclass[../../../include/open-logic-section]{subfiles}

\begin{document}

\olfileid{fol}{int}{mod}

\olsection{Model lan Teori}

Sawise sintaksis lan semantik logika sepisanan ditemtokake, kita bisa
miwiti nliti sipat-sipat !!{structure}s lan gagasan semantis. Kita uga
bisa nemtokake sistem !!{derivation}, banjur nliti sistem kasebut.
Tumrap sawijining himpunan !!{sentence}s, kita bisa takon: !!{structure}s
endi kang ndadekake kabeh !!{sentence}s ing himpunan kasebut bener?
Yen diwenehi himpunan !!{sentence}s~$\Gamma$, !!a{structure}~$\Struct{M}$
kang nyukupi ukara-ukara kasebut diarani \emph{model saka~$\Gamma$}.
Kita bisa miwiti saka~$\Gamma$ lan nyoba nemokake model-modele---kaya
apa rupane? Sepira gedhe utawa cilike? Nanging kita uga bisa miwiti
saka siji !!{structure} utawa sawenehing !!{structure}s banjur takon:
!!{sentence}s endi kang bener ing struktur-struktur kasebut? Apa ana
!!{sentence}s kang \emph{menehi ciri} marang !!{structure}s iki kanthi
pangerten yen ukara-ukara kasebut bener ing struktur-struktur iki, lan
mung ing struktur-struktur iki? Cathetan koreksi sumber OLPL-062:
sumber Inggris nyelehake tembung ``mung'' marang pronomina kanggo
ukara, kang kanthi harfiah malah mbatesi ukara kang bener; konteks teori
model lan metode aksiomatis mbutuhake ukara kasebut bener persis ing
struktur saka kelas kang diwenehi ciri. Pitakon-pitakon kaya mangkono
kalebu wewengkon \emph{teori model}. Pitakon kasebut uga dadi dhasar
kanggo \emph{metode aksiomatis}: njlentrehake sawijining kumpulan
!!{structure}s nganggo himpunan !!{sentence}s, yaiku aksioma-aksioma
sawijining teori. Iki bisa ditindakake amarga pangamatan yen persis
!!{sentence}s kang dadi akibat saka aksioma-aksioma ing logika
sepisanan iku bener ing kabeh model aksioma kasebut.

Minangka tuladha kang prasaja banget, timbangana pratatanan. Sawijining
pratatanan yaiku relasi~$R$ ing sawenehing himpunan~$A$ kang refleksif
lan transitif. Himpunan~$A$ kanthi relasi rong panggonan $R \subseteq A
\times A$ ing himpunan kasebut persis kaya kang dibutuhake kanggo
menehi !!a{structure} tumrap basa logika sepisanan kang mung nduweni
siji simbol relasi rong panggonan~$\Obj P$: kita netepake
$\Domain{M} = A$ lan $\Assign{\Obj P}{M} = R$. Amarga $R$ iku
pratatanan, relasi kasebut refleksif lan transitif, lan kita bisa
nemokake himpunan~$\Gamma$ saka !!{sentence}s logika sepisanan kang
nyatakake iki:
\begin{align*}
  & \lforall[\Obj v_0][\Atom{\Obj P}{\Obj v_0,\Obj v_0}]\\
  & \lforall[\Obj v_0][\lforall[\Obj v_1][\lforall[\Obj v_2][((\Atom{\Obj P}{\Obj v_0,\Obj v_1} \land \Atom{\Obj P}{\Obj v_1,\Obj v_2}) \lif \Atom{\Obj P}{\Obj v_0,\Obj v_2})]]]
\end{align*}
!!{sentence}s kasebut mung simbolisasi saka ``kanggo sembarang~$x$,
$Rxx$'' ($R$ refleksif) lan ``saben $Rxy$ lan $Ryz$, uga $Rxz$''
($R$ transitif). Kita bisa ndeleng yen !!a{structure}~$\Struct{M}$ iku
model saka loro !!{sentence}s~$\Gamma$ iki yen lan mung yen $R$
(yaiku, $\Assign{\Obj P}{M}$) dadi pratatanan ing~$A$ (yaiku,
$\Domain{M}$). Kanthi tembung liya, model-model saka~$\Gamma$ persis
pratatanan. Saben sipat saka kabeh pratatanan kang bisa diandharake ing
basa logika sepisanan kang mung nggunakake~$\Obj P$ minangka
!!{predicate} (kayata refleksivitas lan transitivitas ing ndhuwur),
dadi akibat saka loro !!{sentence}s ing~$\Gamma$, lan kosok balene.
Dadi, apa wae kang bisa kita buktekake ngenani model-model saka~$\Gamma$
wis kita buktekake ngenani kabeh pratatanan.

Tumrap saben teori lan kelas model tartamtu (kayata $\Gamma$ lan kabeh
pratatanan), bakal ana pitakon-pitakon kang narik kawigaten ngenani apa
kang bisa diandharake ing basa logika sepisanan kang gegandhengan, lan
apa kang ora bisa diandharake. Ana sawetara sipat !!{structure}s kang
narik kawigaten tumrap kabeh basa lan kelas model, yaiku sipat kang
magepokan karo gedhene !!{domain}. Tuladhane, kita mesthi bisa
ngandharake yen !!{domain} ngemot persis $n$~!!{element}s, kanggo
sembarang~$n \in \PosInt$. Kita uga bisa ngandharake, kanthi
nggunakake himpunan tanpa wates saka !!{sentence}s, yen !!{domain}
tanpa wates. Nanging, kita ora bisa ngandharake yen domain iku winates,
utawa yen domain iku !!{nonenumerable}. Asil-asiling wates logika
sepisanan iki minangka akibat saka teorema kekompakan lan
L\"owenheim--Skolem.

\end{document}
